En el model clàssic de l'àtom d'hidrogen, l'electró gira al voltant del protó en una òrbita circular de radi $r = 53\ \mathrm{pm}$.

\begin{apartat}{1}
Calculeu el mòdul de la força elèctrica que actua sobre l'electró. Representeu aquesta força en dos punts de l'òrbita amb una separació angular de 90°. Calculeu el mòdul del camp elèctric que crea el protó en un punt de la trajectòria de l'electró.
\end{apartat}

\begin{apartat}{1}
Calculeu l'energia mecànica d'aquest sistema, que consta d'un protó i un electró girant al seu voltant. Expresseu el resultat en eV.
\end{apartat}

\dades{%
  $k = \dfrac{1}{4\pi\varepsilon_0} = 8,99\times 10^{9}\ \mathrm{N\,m^2\,C^{-2}}$.\\
  Càrrega de l'electró, $q_{\text{electró}} = -1,6\times 10^{-19}\ \mathrm{C}$.\\
  Càrrega del protó, $q_{\text{protó}} = 1,6\times 10^{-19}\ \mathrm{C}$.\\
  Massa de l'electró, $m_{\text{electró}} = 9,11\times 10^{-31}\ \mathrm{kg}$.\\
  $1\ \mathrm{eV} = 1,60\times 10^{-19}\ \mathrm{J}$.}
